\section{Experiments}
\label{sec:experiments}

\subsection{Setup}
All experiments run on a single consumer GPU. Synthetic sequences are
generated by our controllable thermal scene model (drifting correlated
background, moving hot/cold objects with soft thermal edges, microbolometer
low-pass, PRNU, NETD noise, 14-bit quantization, AGC drift/jumps). Real
sources are HIT-UAV sample videos~\cite{suo2023hituav} (\SI{7.5}{\hertz},
$640{\times}512$) and the eight continuous sequences of the FLIR ADAS v2
video test set~\cite{fliradas} ($640{\times}512$, 3{,}145 frames total,
COCO labels).

\subsection{E1: interpolation audit}
\label{sec:e1}
We hold out every $k$-th frame from real thermal video, reconstruct it from
the remaining frames, and measure PSNR/SSIM for linear, cubic, and
flow-guided interpolation. Table~\ref{tab:e1} shows that
\emph{linear blending matches or beats flow-guided warping} on all five
HIT-UAV videos at $4\times$ and $8\times$ upsampling. Thermal imagery is
texture-poor, AGC-corrupted, and at \SI{7.5}{\hertz} exhibits large
non-rigid platform motion, violating the brightness-constancy assumption
that flow-guided (and learned, visible-light-trained) interpolators rely
on. We therefore default to linear interpolation, and flag the failure of
visible-light interpolation practice to transfer to thermal as an E1
finding.

\begin{table}[t]
\centering
\caption{Held-out frame reconstruction on HIT-UAV thermal video (PSNR dB).
Linear blending is competitive or best; flow guidance does not transfer.}
\label{tab:e1}
\begin{tabular}{lccc}
\toprule
video & linear & cubic & flow \\
\midrule
\multicolumn{4}{l}{\emph{values filled from experiments/e1\_interp\_results.json}}\\
\bottomrule
\end{tabular}
\end{table}

\subsection{E2: conversion and AGC de-flicker}
\label{sec:e2}
We convert all real sources with our pipeline and measure event statistics
and flicker before/after de-flickering. On synthetic scenes with known
ground truth, pivot-centred AGC correction reduces the static-scene
spurious event count to the \emph{exact} radiometric floor
(Fig.~\ref{fig:agc}), whereas naive two-parameter affine correction
increases the event rate by an order of magnitude by unpinning bulk
pixels.

\subsection{E3: fidelity and parameter sensitivity}
\label{sec:fidelity}
\label{sec:calibration}
We sweep simulator mode $\times$ threshold $\times$ shot noise $\times$
bolometer low-pass on fixed ground-truth sequences and report event-rate
statistics, polarity balance, spatial sparsity, and Edge--Event Alignment
(EEA) as a function of parameters. We further calibrate the
threshold$\leftrightarrow$NETD mapping by measuring event yield inside
known object boxes as object temperature contrast is swept
(Fig.~\ref{fig:calibration}), giving the minimum detectable contrast per
threshold. Results are reported as sensitivity \emph{bands}, not single
points.

\subsection{E4: detection benchmark}
\label{sec:e4}
We train a shared CenterNet-style detector~\cite{zhou2019centernet} on
three input channels---thermal frame, event voxel grid, and early
fusion---plus a recurrent (ConvLSTM) event variant, isolating the input
representation as the only varied factor. An external YOLOv8n anchor is
trained on the FLIR thermal frame train split. We report COCO mAP and
mAP@50 overall and per slice (small/large, slow/fast, low/high contrast,
day/night, AGC-corrupted) on both STED and FLIR-STED.
